\chapter{Den tause kunnskapen er ikkje privat}

\begin{motto}
At noko ikkje let seg seie med ein gong,\\ tyder ikkje at det ikkje let seg lære.
\end{motto}

Den mest øydeleggjande skuldinga i \textsc{Grunnlaust} er ikkje at design manglar metode eller etikk. Det er at designkunnskapen er privat — taus heile vegen ned, låst inne i den einskilde utøvaren si magekjensle, umogleg å gjere eksplisitt, umogleg å prøve, umogleg å overføre. Om dette stemmer, hjelper det lite at design har eit eige studieobjekt og eit eige tidsskrift; eit fag der kvar utøvar veit noko ingen andre kan etterprøve, er ikkje eit fag, men ei samling private intuisjonar. Dette er den alvorlegaste innvendinga, og dette kapitlet svarar på henne direkte. Påstanden er at den tause kunnskapen verken er privat eller særeigen for design — at han er ekte, eksaminerbar, for det meste kollektiv, og overførbar gjennom den same opplæringa som ber all profesjonell kompetanse. Og at den som peikar ut design som unikt fundamentlaust fordi kunnskapen er taus, beviser for mykje: same argumentet ville felle fysikken.

\section{Schön: refleksjon-i-handling er ekte kunnskap}

Donald Schön bygde heile forfattarskapen sin kring eit einaste, presist poeng: at den verkelege kompetansen til ein profesjonell ligg ikkje i å bruke teori på praksis, men i ein eigen form for tenking som skjer \emph{i} handlinga \autocite{schon1983}. Han kalla det rådande biletet \emph{teknisk rasjonalitet} — ideen om at den dugande utøvaren fyrst lærer ein vitskapleg teori og så bruker han på problem som ventar ferdig formulerte. Gjennom nære studiar av korleis arkitektar, ingeniørar, legar og psykoterapeutar faktisk arbeider, viste Schön at dette ikkje skildrar nokon profesjon. Den verkelege kompetansen er ein \emph{refleksjon-i-handling}: utøvaren møter ein situasjon som ikkje passar nokon ferdig kategori, rammar han på ein måte, prøver eit grep, les korleis situasjonen svarar, og rammar på nytt i lys av svaret. Schön kalla det ein refleksiv samtale med materialet i situasjonen \autocite{schon1992designing}.

Schön gav denne tenkinga ein presis indre struktur, og det er verdt å hente fram, for det er strukturen som gjer at ho lèt seg prøve. Utøvaren ber med seg det Schön kallar eit repertoar — eit lager av tidlegare tilfelle, bilete og grep — og møter den nye situasjonen ved å sjå han \emph{som} eit av dei kjende tilfella, samstundes som han er var for kvar den nye situasjonen avvik. Han gjer det Schön kallar eit ramme-eksperiment: han set ein ramme på situasjonen, handlar ut frå henne, og lèt situasjonen «svare tilbake» — materialet, teikninga, brukaren, konstruksjonen gjer motstand eller gjev etter på måtar utøvaren ikkje fullt ut styrer. Denne tilbaketalen er det som disiplinerer prosessen. Utøvaren kan ikkje berre bestemme at grepet var godt; han må lese korleis situasjonen tek imot det, og endre seg etter svaret. Dette er strengt i ordets eigentlege tyding: prosessen er underlagd ein ytre kontroll han ikkje rår over åleine.

Det avgjerande er at dette ikkje er fråvær av kunnskap. Det er kunnskap av eit anna slag — og, det Schön viste tydelegast, kunnskap som kan vere god eller dårleg, dugande eller udugande. Ein utøvar kan ramme problemet feil, lese responsen feil, halde fast ved eit grep som situasjonen alt har avvist. Ein annan kan ramme det treffande og lese det fint. Skilnaden mellom dei to er ikkje smak. Han er kompetanse, og han kan vurderast av andre kompetente utøvarar. Dette er heile poenget mot \textsc{Grunnlaust}: refleksjon-i-handling er ikkje ei bortforklaring av at design ikkje kan grunngje seg. Det er ein skildring av korleis design \emph{faktisk} grunngjev seg — ikkje ved å utleie konklusjonar frå premissar, men ved ein disiplinert dialog med ein situasjon som svarar tilbake. At grunngjevinga ikkje har deduksjonens form, gjer henne ikkje privat. Han har ei anna form, og forma er studert.

\section{Studioet: korleis det taus blir rigorøst lært}

Om den tause kunnskapen var verkeleg privat, ville det vere umogleg å lære han vidare. Men design \emph{blir} lært vidare, kvar dag, på ein måte som er så velprøvd at han har vore stabil i over hundre år: gjennom atelieret. I \emph{Educating the Reflective Practitioner} gjorde Schön designstudioet til paradigmet på korleis ein indeterminert, skjønsbasert praksis kan lærast strengt \autocite{schon1987educating}. Studioet er ikkje ein stad der ein overfører reglar. Det er ein stad der studenten gjer noko, viser fram det ho har gjort, og får det lese av ein meister og av jamaldringar i ein offentleg kritikk. Læraren peikar, demonstrerer, omformulerer; studenten prøver på nytt. Over tid lærer studenten ikkje ein katalog av reglar, men ein måte å sjå og handle på — ein måte å ramme problem, kjenne att det lovande grepet, og dømme sitt eige arbeid.

Det viktige er kor mykje av dette som faktisk er offentleg. Kritikken i studioet er ikkje privat; han skjer framfor heile gruppa, med eksplisitte grunngjevingar som andre kan motseie. Når ein erfaren kritikar seier at ein plan ikkje held, peikar ho på \emph{kvifor} — på kva relasjon som ikkje stemmer, kva omsyn som er gløymt, kva grep som ville ha løyst det. Grunngjevinga er ikkje alltid uttømmande, men ho er ekte, og ho lèt seg vurdere. Det same skjer i kvar profesjonsutdanning der skjøn er sentralt: legen lærer klinisk skjøn ved sengekanten, advokaten lærer juridisk dømmekraft i møtet med verkelege saker. At noko blir lært gjennom demonstrasjon og kritikk heller enn gjennom lærebok, gjer det ikkje uoverførbart. Det gjer det overført \emph{på den måten praktisk kunnskap alltid har vore overført}. Studioet er ikkje eit teikn på at design manglar ein metode for å føre kunnskapen vidare. Det \emph{er} metoden, og det er ein god ein.

\section{Polanyi: alt vi veit har ein taus botn}

Her kjem det avgjerande filosofiske svaret, og det snur skuldinga på hovudet. \textsc{Grunnlaust} skriv som om taus kunnskap var eit problem særeige for design — eit teikn på at faget aldri nådde fram til eksplisitt, etterprøvbar viten. Men Michael Polanyi viste at det motsette er sant: \emph{all} kunnskap kviler på ein taus dimensjon \autocite{polanyi1966tacit}. Polanyi sin berømte formel er at vi veit meir enn vi kan seie. Vi kjenner att eit ansikt blant tusen utan å kunne seie kva trekk vi går etter. Vi sit på sykkelen utan å kunne formulere fysikken som held oss oppreist. Og — dette er Polanyi sitt skarpaste poeng — fysikaren les måleinstrumenta sine, vurderer kva som tel som eit reint forsøk, og kjenner att eit lovande spor, alt saman ved hjelp av eit innøvd skjøn som han ikkje kan gjere fullt eksplisitt. Den vitskaplege kunnskapen, som \textsc{Grunnlaust} held opp som mønster på det eksplisitte, har sjølv ein taus botn ho ikkje kan kvitte seg med.

Dette beviser for mykje for kritikaren. Om taus kunnskap diskvalifiserer eit fag frå å ha eit fundament, så har ingen fag eit fundament — ikkje fysikken, ikkje matematikken, ikkje noko. Argumentet, ført til endes, et opp seg sjølv. Og då har kritikaren to val. Anten godtek han at taus kunnskap er foreinleg med eit fundament — i kva tilfelle skuldinga mot design fell, av di designkunnskapen sin tause dimensjon ikkje skil han frå nokon annan disiplin. Eller han held fast ved at taus kunnskap utelukkar eit fundament — i kva tilfelle han har felt heile vitskapen saman med design. Han kan ikkje ha det begge vegar. Det er ingen veg som rammar design utan òg å ramme det faget kritikaren held design opp mot. Polanyi gjer ikkje design til eit unntak. Han viser at det aldri var eit unntak.

\section{Ryle og Collins: knowing-how er kunnskap, og han er felles}

To fleire steg trengst for å lukke argumentet. Det fyrste er Gilbert Ryle sitt skilje mellom å vite \emph{at} og å vite \emph{korleis} \autocite{ryle1949mind}. Ryle gjorde opp med det han kalla den intellektualistiske legenda: ideen om at all ekte kunnskap eigenleg er kunnskap om sanne påstandar, og at praktisk dugleik berre er anvend teori. Han viste at dette er feil. Å vite korleis ein gjer noko — korleis ein argumenterer godt, spelar sjakk, set ein setning — er ikkje å ha ein indre regelbok ein konsulterer. Det er ein kunnskap i seg sjølv, fullverdig, som viser seg i å gjere tingen godt og ikkje i å kunne seie korleis. Klokskap er ikkje teori om klokskap. Når \textsc{Grunnlaust} avviser designaren sitt skjøn fordi det ikkje kjem i påstandsform, gjer det nett feilen Ryle avslørte: det reknar berre å-vite-at som kunnskap, og ser ikkje at å-vite-korleis er like ekte. Designkunnskap er framifrå å-vite-korleis. At han ikkje let seg skrive om til ei liste påstandar, er ikkje ein mangel ved han. Det er kva slags kunnskap han er.

Det andre steget, og det som direkte slår beina under «privat heile vegen ned», kjem frå Harry Collins \autocite{collins2010tacit}. Collins har gjort det grundigaste arbeidet vi har på taus kunnskap, og hovudfunnet hans er at det aller meste av tause kunnskap ikkje er privat i det heile. Han skil mellom fleire slag. Noko er taust berre i praksis — det kunne skrivast ned, men ingen har gjort det enno. Noko er taust av di det krev ein kropp — som balanse, som ikkje kan overførast i ord men berre i øving. Og det viktigaste slaget, som Collins kallar \emph{kollektiv} taus kunnskap, finst ikkje i hovudet til den einskilde i det heile; det bur i eit praksisfellesskap, og blir tileigna ved å delta i det. Å lære eit fags tause kunnskap er å bli sosialisert inn i ein måte å gjere ting på som gruppa ber i lag. Dette er sjølve motsetnaden til privat. Kunnskapen er felles eigedom, overført ved deltaking, og difor i prinsippet tilgjengeleg for kven som helst som går inn i fellesskapet og lærer.

Med dette er biletet i \textsc{Grunnlaust} snudd heilt om. Designkunnskapen er ikkje ein einsleg sin uutgrunnelege intuisjon. Han er ein kollektiv, innøvd, overførbar kompetanse — ekte som å-vite-korleis, felles som eit handverk, og ikkje meir taus enn kunnskapen i kvart anna fag der dugleik tel. At han for ein stor del lèt seg seie berre delvis, gjer han ikkje privat. Det gjer han menneskeleg, slik nær sagt all viktig kunnskap er. Eg skal vere ærleg om kvar grensa går: det følgjer ikkje av dette at \emph{kvart} ord i designaren sitt dommarvokabular tyder akkurat det same for alle, og spørsmålet om kor presist fagspråket er kjem boka tilbake til. Men hovudsaka står. Den tause kunnskapen er ikkje privat, og difor fell den alvorlegaste skuldinga.

\vignett

Står att eitt spørsmål: sjølv om designkunnskapen er felles og overførbar, har han ein eigen \emph{logikk}? Finst det ei slutningsform som er designens eiga, like formell som deduksjon og induksjon? Det er det neste kapittelet svarar ja på.

\vidarelesing{\fullcite{schon1983}; \fullcite{schon1987educating}; \fullcite{schon1992designing}; \fullcite{polanyi1966tacit}; \fullcite{ryle1949mind}; \fullcite{collins2010tacit}.}
